\documentclass[11pt,a4paper]{article}

\usepackage[utf8]{inputenc}
\usepackage[T1]{fontenc}
\usepackage[turkish]{babel}
\usepackage[margin=2.4cm]{geometry}
\setlength{\headheight}{14pt}
\usepackage{graphicx}
\usepackage{float}
\usepackage{booktabs}
\usepackage{tabularx}
\usepackage{amsmath,amssymb}
\usepackage{gensymb}
\usepackage[hidelinks]{hyperref}
\usepackage{fancyhdr}
\usepackage{enumitem}
\usepackage{xcolor}
\usepackage{cite}

\definecolor{bugred}{rgb}{0.75,0.1,0.1}
\definecolor{solgreen}{rgb}{0.1,0.5,0.1}

\pagestyle{fancy}
\fancyhf{}
\fancyhead[L]{\small Alt Ekstremite ROM -- Sens\"or Raporu}
\fancyhead[R]{\small \thepage}
\renewcommand{\headrulewidth}{0.4pt}

% ================================================================
\begin{document}

% -- Baslik --
\begin{center}
    {\LARGE\bfseries HTC VIVE Ultimate Tracker ile\\[4pt]
    Alt Ekstremite Eklem ROM \"Ol\c{c}\"um\"u:\\[4pt]
    Sens\"or Yerle\c{s}imi ve Kalibrasyon Raporu}\\[12pt]
    {\normalsize Mart 2026}
\end{center}

\vspace{6pt}

\begin{abstract}
Bu rapor, \"u\c{c} adet HTC VIVE Ultimate Tracker kullan{\i}larak sa\u{g} alt
ekstremitenin (kal\c{c}a ve diz eklemi) hareket a\c{c}{\i}kl{\i}\u{g}{\i}n{\i}n (ROM) ger\c{c}ek
zamanl{\i} \"ol\c{c}\"ulmesine y\"onelik bir sistemin sens\"or yerle\c{s}im kararlar{\i}n{\i},
kalibrasyon y\"ontemini ve tespit edilen yaz{\i}l{\i}m sorunlar{\i}n{\i}
sunmaktad{\i}r. Tracker'lar pelvis, femur ve tibia segmentlerine
ba\u{g}lanarak forward kinematics (FK) ile eklem a\c{c}{\i}lar{\i}
hesaplanmaktad{\i}r. Kalibrasyon, statik referans poz y\"ontemiyle
ger\c{c}ekle\c{s}tirilmi\c{s}tir. Test s\"urecinde iskelet hiyerar\c{s}isinden
kaynaklanan iki hata tespit edilmi\c{s} ve k\"ok neden analizi
yap{\i}lm{\i}\c{s}t{\i}r.
\end{abstract}

\tableofcontents
\vspace{8pt}

% ================================================================
\section{Giri\c{s}}
\label{sec:giris}

Eklem hareket a\c{c}{\i}kl{\i}\u{g}{\i} (ROM) \"ol\c{c}\"um\"u, ortopedik de\u{g}erlendirme ve
rehabilitasyon s\"urecinde temel bir klinik parametredir~\cite{norkin2016}.
Geleneksel olarak universal gonyometre ile yap{\i}lan bu \"ol\c{c}\"um, g\"ozlemciler
aras{\i} de\u{g}i\c{s}kenlik ve statik \"ol\c{c}\"um s{\i}n{\i}rl{\i}l{\i}\u{g}{\i} ta\c{s}{\i}r~\cite{gajdosik1987}.
Giyilebilir ataletsel \"ol\c{c}\"um birimleri (IMU) ve optik tabanl{\i} takip
sistemleri, bu s{\i}n{\i}rl{\i}l{\i}klar{\i} a\c{s}mak i\c{c}in alternatif olarak
kullan{\i}lmaktad{\i}r~\cite{cuesta2010,seel2014}.

HTC VIVE Ultimate Tracker, iki geni\c{s} g\"or\"u\c{s} a\c{c}{\i}l{\i} kamera ve bilgisayarl{\i}
g\"or\"u (computer vision) algoritmalar{\i} ile 6-serbestlik dereceli (6DoF)
inside-out takip sa\u{g}layan kompakt bir sens\"ord\"ur
(77$\times$58.6$\times$27.3\,mm, 94\,g)~\cite{htcvive2024}. Baz istasyonu
gerektirmeden \c{c}al{\i}\c{s}mas{\i}, klinik ortamlarda pratik bir avantaj
sunmaktad{\i}r~\cite{niehorster2017}.

Bu \c{c}al{\i}\c{s}mada, \"u\c{c} adet VIVE Ultimate Tracker kullan{\i}larak sa\u{g} alt
ekstremitenin kal\c{c}a ve diz fleksiyon/ekstansiyon a\c{c}{\i}lar{\i}
\"ol\c{c}\"ulm\"u\c{s}t\"ur. Sistem Unity~6 ve VIVE OpenXR Plugin~2.5.1 altyap{\i}s{\i}
\"uzerinde \c{c}al{\i}\c{s}maktad{\i}r.

% ================================================================
\section{VIVE Ultimate Tracker: \c{C}al{\i}\c{s}ma Prensibi}
\label{sec:tracker}

VIVE Ultimate Tracker, \textbf{inside-out} takip mimarisine sahiptir;
yani konum ve y\"onelim bilgisini kendi \"uzerindeki sens\"orlerden
t\"uretir, d{\i}\c{s} bir baz istasyonuna ihtiya\c{c} duymaz~\cite{htcvive2024}.

\textbf{Optik takip (pozisyon + y\"onelim):}
Tracker \"uzerinde iki adet geni\c{s} g\"or\"u\c{s} a\c{c}{\i}l{\i} (wide-FOV) kamera
bulunur. Bu kameralar \c{c}evredeki g\"orsel \"ozellikleri (kenarlar,
k\"o\c{s}eler, doku farklar{\i}) ger\c{c}ek zamanl{\i} olarak alg{\i}lar. Bilgisayarl{\i}
g\"or\"u (SLAM -- Simultaneous Localization and Mapping) algoritmalar{\i},
ard{\i}\c{s}{\i}k kareler aras{\i}ndaki g\"orsel yer de\u{g}i\c{s}tirmeyi analiz ederek
tracker'{\i}n 6DoF pozunu (3 eksen pozisyon + 3 eksen y\"onelim)
hesaplar~\cite{cadena2016}.

\textbf{Ataletsel \"ol\c{c}\"um birimi (IMU):}
Dahili IMU (ivme\"ol\c{c}er + jiroskop), y\"uksek frekansl{\i} rotasyon
de\u{g}i\c{s}imlerini yakalar ve optik takibin zaman \"oz\"un\"url\"u\u{g}\"un\"u
art{\i}r{\i}r. Kameralar ortam{\i}n g\"orsel haritasna dayanarak mutlak
pozisyon sa\u{g}larken, IMU k{\i}sa s\"ureli h{\i}zl{\i} hareketlerde
s\"ur\"ukleni\c{s}i (drift) telafi eder. \.Iki sens\"or tipi bir f\"uzyon
filtresiyle birle\c{s}tirilerek d\"u\c{s}\"uk gecikmeli, kararl{\i} bir poz
\"uretilir~\cite{seel2014}.

\textbf{Kablosuz ba\u{g}lant{\i}:}
Tracker, 2.4\,GHz tescilli protokol \"uzerinden VIVE Wireless Dongle'a
ba\u{g}lan{\i}r. Tek bir dongle, e\c{s} zamanl{\i} olarak 5 adede kadar
tracker'{\i} destekler. Ba\u{g}lant{\i} gecikme s\"uresi d\"u\c{s}\"uk tutularak
ger\c{c}ek zamanl{\i} hareket takibi sa\u{g}lan{\i}r.

Tablo~\ref{tab:tracker_bilesenler} tracker bile\c{s}enlerini \"ozetler.

\begin{table}[H]
\centering
\caption{VIVE Ultimate Tracker bile\c{s}enleri ve i\c{s}levleri}
\label{tab:tracker_bilesenler}
\renewcommand{\arraystretch}{1.2}
\begin{tabularx}{\textwidth}{lXl}
\toprule
\textbf{Bile\c{s}en} & \textbf{\.I\c{s}lev} & \textbf{\c{C}{\i}k{\i}\c{s}} \\
\midrule
Wide-FOV kamera $\times$2 & SLAM ile mutlak pozisyon ve y\"onelim & 6DoF poz \\
IMU (ivme\"ol\c{c}er + jiroskop) & Y\"uksek frekansl{\i} rotasyon takibi, drift telafisi & A\c{c}{\i}sal h{\i}z, ivme \\
Yapay zek\^a motoru & G\"or\"u + IMU f\"uzyonu, poz tahmini & Filtrelenmi\c{s} poz \\
2.4\,GHz radyo & HMD'ye kablosuz veri iletimi & Poz verisi \\
\bottomrule
\end{tabularx}
\end{table}

% ================================================================
\section{Sens\"or Yerle\c{s}imi}
\label{sec:yerlesim}

\subsection{Yerle\c{s}im Kriterleri}

Sens\"or konumlar{\i} belirlenirken a\c{s}a\u{g}{\i}daki \"ol\c{c}\"utler dikkate
al{\i}nm{\i}\c{s}t{\i}r~\cite{cuesta2010,niswander2020}:

\begin{enumerate}[nosep]
    \item \textbf{Yumu\c{s}ak doku artefakt{\i}n{\i}n minimizasyonu:} Sens\"or, kemik
          y\"uzeyine yak{\i}n b\"olgelere yerle\c{s}tirilerek kas kas{\i}lmas{\i} ve
          deri kaymas{\i}ndan kaynaklanan hata azalt{\i}l{\i}r~\cite{leardini2005}.
    \item \textbf{Segment temsil yeterlili\u{g}i:} Her sens\"or, ba\u{g}l{\i} oldu\u{g}u
          rijit kemik segmentinin rotasyonunu do\u{g}ru temsil etmelidir.
    \item \textbf{Hareket s{\i}ras{\i}nda stabilite:} Sens\"or, dinamik hareketlerde
          (y\"ur\"ume, \c{c}\"omelme) kaymamal{\i}d{\i}r.
    \item \textbf{Kom\c{s}u sens\"orlerle \c{c}ak{\i}\c{s}mama:} Tracker'lar birbirine
          fiziksel olarak temas etmemelidir.
\end{enumerate}

\subsection{Nihai Sens\"or Konumlar{\i}}

\begin{table}[H]
\centering
\caption{Tracker yerle\c{s}im konumlar{\i} ve anatomik referanslar{\i}}
\label{tab:konumlar}
\renewcommand{\arraystretch}{1.25}
\begin{tabularx}{\textwidth}{llXl}
\toprule
\textbf{Tracker} & \textbf{Segment} & \textbf{Anatomik Konum} & \textbf{Sabitleme} \\
\midrule
0 -- Pelvis & G\"ovde & Umbilikus hizas{\i}, anterior abdomen, orta hat & Elastik bant \\
1 -- Femur  & \"Ust bacak & Patella proksimali ($\sim$5--10\,cm), femur distal 1/3 anterolateral & Elastik bant \\
2 -- Tibia  & Alt bacak & Lateral malleolus proksimali ($\sim$5--8\,cm), tibia distal 1/3 anterolateral & Elastik bant \\
\bottomrule
\end{tabularx}
\end{table}

Pelvis tracker g\"obek b\"olgesine yerle\c{s}tirilmi\c{s}tir \c{c}\"unk\"u Mixamo iskelet
modelinde \texttt{pelvisBone} olarak kullan{\i}lan kemik, g\"ovde
rotasyonunu temsil etmektedir. Femur tracker'{\i}, diz kapa\u{g}{\i}n{\i}n
hemen proksimaline konumland{\i}r{\i}lm{\i}\c{s}t{\i}r; bu b\"olge anteriorda kemik
y\"uzeyine yak{\i}nd{\i}r ve yumu\c{s}ak doku artefakt{\i} d\"u\c{s}\"ukt\"ur~\cite{leardini2005}.
Tibia tracker'{\i}, tibian{\i}n distal anterolateral y\"uzeyine
yerle\c{s}tirilmi\c{s}tir; bu b\"olgede subk\"utan yumu\c{s}ak doku kal{\i}nl{\i}\u{g}{\i}
v\"ucuttaki en d\"u\c{s}\"uk de\u{g}erler aras{\i}ndad{\i}r.

% ================================================================
\section{Kalibrasyon Y\"ontemi}
\label{sec:kalibrasyon}

\subsection{Statik Poz Kalibrasyonu}

Giyilebilir sens\"orlerle eklem a\c{c}{\i}s{\i} \"ol\c{c}\"um\"unde, sens\"or y\"onelimi ile
anatomik segment y\"onelimi aras{\i}ndaki fark{\i}n telafi edilmesi
gerekir. Bu \c{c}al{\i}\c{s}mada \textbf{statik referans poz} (anatomical
calibration) y\"ontemi kullan{\i}lm{\i}\c{s}t{\i}r~\cite{favre2006,palermo2014}. Bu
y\"ontemde kullan{\i}c{\i}, bilinen bir anatomik pozisyonda (dik duru\c{s},
bacaklar d\"uz) hareketsiz durur ve sens\"or-kemik y\"onelim fark{\i}
kaydedilir.

\subsection{Matematiksel Form\"ulasyon}

\subsubsection{Sens\"or-Kemik Offset}
Kalibrasyon an{\i}nda, her tracker ($T$) ile kar\c{s}{\i}l{\i}k gelen kemik ($B$)
aras{\i}ndaki rotasyon fark{\i} hesaplan{\i}r:

\begin{equation}
    Q_{\text{offset}} = Q_{T}^{-1} \cdot Q_{B}
    \label{eq:offset}
\end{equation}

\c{C}al{\i}\c{s}ma zaman{\i}nda kemi\u{g}in d\"unya rotasyonu \c{s}\"oyle elde edilir:

\begin{equation}
    Q_{B}^{t} = Q_{T}^{t} \cdot Q_{\text{offset}}
    \label{eq:runtime}
\end{equation}

\subsubsection{Hiyerar\c{s}ik Lokal Rotasyon (FK)}
Eklem a\c{c}{\i}lar{\i}, ebeveyn-\c{c}ocuk kemik ili\c{s}kisinden t\"uretilir. Forward
kinematics yakla\c{s}{\i}m{\i}yla~\cite{craig2005}:

\begin{align}
    Q_{\text{hip}}^{\text{local}} &= Q_{\text{pelvis}}^{-1} \cdot Q_{\text{thigh}} \label{eq:hip} \\
    Q_{\text{knee}}^{\text{local}} &= Q_{\text{thigh}}^{-1} \cdot Q_{\text{shin}} \label{eq:knee}
\end{align}

Kalibrasyon pozundaki lokal rotasyonlar ($Q^{\text{calib}}$) referans
al{\i}narak delta rotasyon hesaplan{\i}r ve modelin ba\c{s}lang{\i}\c{c} yerel
rotasyonuna uygulan{\i}r:

\begin{equation}
    Q_{\text{joint}}^{\text{model}} = Q_{\text{joint}}^{\text{initial}} \cdot
    \left(Q_{\text{joint}}^{\text{calib}}\right)^{-1} \cdot Q_{\text{joint}}^{\text{current}}
    \label{eq:delta}
\end{equation}

\subsubsection{Klinik A\c{c}{\i} Hesab{\i} (Gonyometri)}
Kal\c{c}a fleksiyon a\c{c}{\i}s{\i}, \"ust bacak y\"on vekt\"or\"un\"un pelvis \c{c}er\c{c}evesinde
sagittal d\"uzleme (Y-Z) projeksiyonuyla hesaplan{\i}r~\cite{norkin2016}:

\begin{equation}
    \theta_{\text{hip}} = -\;\text{atan2}\!\left(v_z^{\text{proj}},\; -v_y^{\text{proj}}\right)
    \label{eq:hip_angle}
\end{equation}

Diz fleksiyon a\c{c}{\i}s{\i}, alt bacak y\"on\"un\"un \"ust bacak \c{c}er\c{c}evesinde
benzer projeksiyonuyla elde edilir:

\begin{equation}
    \theta_{\text{knee}} = \text{atan2}\!\left(v_z,\; -v_y\right)
    \label{eq:knee_angle}
\end{equation}

Pozitif de\u{g}erler fleksiyonu, negatif de\u{g}erler ekstansiyonu ifade eder.

\subsection{Kalibrasyon Prosed\"ur\"u}

\begin{enumerate}[nosep]
    \item Tracker'lar a\c{c}{\i}l{\i}r ve HMD ile e\c{s}le\c{s}tirilir.
    \item Tracker'lar Tablo~\ref{tab:konumlar}'daki konumlara elastik
          bantlarla sabitlenir.
    \item Unity uygulamas{\i} ba\c{s}lat{\i}l{\i}r; \texttt{LegTrackerAutoAssigner}
          aktif tracker'lar{\i} otomatik tespit eder.
    \item Kullan{\i}c{\i} dik durur (bacaklar d\"uz, g\"ovde dik).
    \item Sa\u{g} el kontrolc\"us\"unde A/B butonuna 2\,s bas{\i}l{\i} tutulur
          (veya klavyede \texttt{L} tu\c{s}u).
    \item Ekranda ``Kalibrasyon tamamland{\i}'' mesaj{\i} do\u{g}rulan{\i}r;
          kal\c{c}a ve diz a\c{c}{\i}lar{\i} $0\degree$ g\"ostermelidir.
\end{enumerate}

% ================================================================
\section{Tespit Edilen Hatalar ve K\"ok Neden Analizi}
\label{sec:hatalar}

Test s\"urecinde iki hata g\"ozlemlenmi\c{s}tir. Her ikisinin de \textbf{ayn{\i}
k\"ok nedenden} kaynakland{\i}\u{g}{\i} belirlenmi\c{s}tir.

\subsection{Hata Tan{\i}mlar{\i}}

\begin{itemize}[nosep]
    \item \textcolor{bugred}{\textbf{Hata 1 -- G\"ovde e\u{g}ildi\u{g}inde bacaklar havaya kalk{\i}yor:}}
    Bacaklar d\"uz dururken g\"ovde \"one e\u{g}ildi\u{g}inde, modeldeki bacak
    yukar{\i} do\u{g}ru d\"on\"uyor.
    \item \textcolor{bugred}{\textbf{Hata 2 -- Aksiyel d\"on\"u\c{s}te bacaklar d\"onm\"uyor:}}
    Kullan{\i}c{\i} sa\u{g}a/sola d\"ond\"u\u{g}\"unde g\"ovde d\"oner fakat bacaklar eski
    y\"onlerinde kal{\i}r; bacak kald{\i}r{\i}ld{\i}\u{g}{\i}nda yanl{\i}\c{s} y\"one gider.
\end{itemize}

\subsection{K\"ok Neden: Kemik Hiyerar\c{s}isi Uyu\c{s}mazl{\i}\u{g}{\i}}

Mixamo iskelet yap{\i}s{\i}nda hiyerar\c{s}i \c{s}u \c{s}ekildedir:

\begin{center}
\texttt{Hips} $\rightarrow$ \{\texttt{Spine},\; \texttt{LeftUpLeg},\;
\texttt{RightUpLeg}\}
\end{center}

\texttt{Spine} ve \texttt{RightUpLeg}, \texttt{Hips}'in \textbf{karde\c{s}
(sibling)} \c{c}ocuklar{\i}d{\i}r. Mevcut sistemde \texttt{pelvisBone} olarak
\texttt{mixamorig1:Spine} atanm{\i}\c{s}t{\i}r. Solver, pelvisBone'u
d\"ond\"urd\"u\u{g}\"unde bu rotasyon sadece Spine alt a\u{g}ac{\i}n{\i} etkiler;
karde\c{s} olan RightUpLeg \textbf{etkilenmez}.

\textbf{Hata 1 mekani\u{g}i:} Pelvis tracker \"one e\u{g}ildi\u{g}inde
$Q_{\text{pelvis}}$ de\u{g}i\c{s}ir fakat $Q_{\text{thigh}}$ sabit kal{\i}r.
Denklem~\ref{eq:hip}'e g\"ore $Q_{\text{hip}}^{\text{local}}$ negatif delta
\"uretir; bu, baca\u{g}{\i}n yukar{\i} kalkmas{\i}na neden olur.

\textbf{Hata 2 mekani\u{g}i:} Pelvis tracker yaw d\"on\"u\c{s}\"u yapt{\i}\u{g}{\i}nda
Spine d\"oner fakat Hips (bacaklar{\i}n ger\c{c}ek ebeveyni) d\"onmez;
bacaklar d\"unya koordinat{\i}nda sabit kal{\i}r.

\subsection{\c{C}\"oz\"um \"Onerisi}

\texttt{pelvisBone} atamasinin \texttt{mixamorig1:Spine} yerine
\textcolor{solgreen}{\texttt{mixamorig1:Hips}} olarak de\u{g}i\c{s}tirilmesi
her iki hatay{\i} da \c{c}\"ozer. Hips, t\"um alt hiyerar\c{s}inin (Spine +
bacaklar) ger\c{c}ek ebeveynidir; d\"ond\"ur\"uld\"u\u{g}\"unde t\"um \c{c}ocuk kemikler
birlikte d\"oner.

% ================================================================
\section{Video G\"osterimi}
\label{sec:video}

Sistemin kurulumu, kalibrasyonu ve hata g\"osterimini i\c{c}eren bir video
haz{\i}rlanm{\i}\c{s}t{\i}r.

\textbf{Video i\c{c}eri\u{g}i:} (1)~Tracker'lar{\i}n v\"ucuda ba\u{g}lanmas{\i},
(2)~Unity uygulamas{\i}n{\i}n ba\c{s}lat{\i}lmas{\i} ve tracker alg{\i}lama,
(3)~Dik duru\c{s}ta kalibrasyon (A/B butonu, 2\,s),
(4)~Kal\c{c}a ve diz fleksiyon testleri,
(5)~Hata 1 ve 2'nin canl{\i} g\"osterimi.

% ================================================================
\section{Sonu\c{c} ve Gelecek \c{C}al{\i}\c{s}malar}
\label{sec:sonuc}

VIVE Ultimate Tracker'{\i}n 6DoF inside-out takip yetene\u{g}i, baz
istasyonu gerektirmeden klinik ortamda alt ekstremite ROM \"ol\c{c}\"um\"u
i\c{c}in uygulanabilir bir \c{c}\"oz\"um sunmaktad{\i}r. Statik poz
kalibrasyonu, farkl{\i} kullan{\i}c{\i}lara ve montaj a\c{c}{\i}lar{\i}na uyum
sa\u{g}lamaktad{\i}r. Tespit edilen iskelet hiyerar\c{s}isi hatas{\i}
(\texttt{Spine} $\neq$ \texttt{Hips}), forward kinematics
sistemlerinde ebeveyn-\c{c}ocuk ili\c{s}kisinin kritik \"onemini ortaya
koymaktad{\i}r.

% ================================================================
\begin{thebibliography}{99}

\bibitem{norkin2016}
C.~C.~Norkin and D.~J.~White, \textit{Measurement of Joint Motion: A Guide
to Goniometry}, 5th~ed. Philadelphia: F.A.~Davis Company, 2016.
\url{https://books.google.com/books?id=TSluDQAAQBAJ}

\bibitem{gajdosik1987}
R.~L.~Gajdosik and R.~W.~Bohannon, ``Clinical measurement of range of
motion: Review of goniometry emphasizing reliability and validity,''
\textit{Physical Therapy}, vol.~67, no.~12, pp.~1867--1872, 1987.
\url{https://doi.org/10.1093/ptj/67.12.1867}

\bibitem{cuesta2010}
A.~Cuesta-Vargas, A.~Gal\'an-Mercant, and J.~M.~Williams, ``The use of
inertial sensors system for human motion analysis,''
\textit{Physical Therapy Reviews}, vol.~15, no.~6, pp.~462--473, 2010.
\url{https://doi.org/10.1179/1743288X11Y.0000000006}

\bibitem{seel2014}
T.~Seel, J.~Raisch, and T.~Schauer, ``IMU-based joint angle measurement
for gait analysis,'' \textit{Sensors}, vol.~14, no.~4, pp.~6891--6909,
2014.
\url{https://doi.org/10.3390/s140406891}

\bibitem{htcvive2024}
HTC Corporation, ``VIVE Ultimate Tracker -- Specifications,'' 2024.
[Online]. Available:
\url{https://www.vive.com/us/accessory/vive-ultimate-tracker/}.

\bibitem{niehorster2017}
D.~C.~Niehorster, L.~Li, and M.~Lappe, ``The accuracy and precision of
position and orientation tracking in the HTC Vive virtual reality system
for scientific research,'' \textit{i-Perception}, vol.~8, no.~3, 2017.
\url{https://doi.org/10.1177/2041669517708205}

\bibitem{leardini2005}
A.~Leardini, L.~Chiari, U.~Della~Croce, and A.~Cappozzo, ``Human
movement analysis using stereophotogrammetry -- Part~3: Soft tissue
artifact assessment and compensation,'' \textit{Gait \& Posture}, vol.~21,
no.~2, pp.~212--225, 2005.
\url{https://doi.org/10.1016/j.gaitpost.2004.05.002}

\bibitem{niswander2020}
W.~Niswander, W.~Wang, and K.~Kontson, ``Optimization of IMU sensor
placement for the measurement of lower limb joint kinematics,''
\textit{Sensors}, vol.~20, no.~21, p.~5993, 2020.
\url{https://doi.org/10.3390/s20215993}

\bibitem{favre2006}
J.~Favre, B.~M.~Jolles, O.~Siegrist, and K.~Aminian, ``Quaternion-based
fusion of gyroscopes and accelerometers to improve 3D angle
measurement,'' \textit{Electronics Letters}, vol.~42, no.~11,
pp.~612--614, 2006.
\url{https://doi.org/10.1049/el:20060124}

\bibitem{palermo2014}
E.~Palermo, S.~Rossi, F.~Marini, F.~Patanè, and P.~Cappa,
``Experimental evaluation of accuracy and repeatability of a novel
body-to-sensor calibration procedure for inertial sensor-based gait
analysis,'' \textit{Measurement}, vol.~52, pp.~145--155, 2014.
\url{https://doi.org/10.1016/j.measurement.2014.03.004}

\bibitem{craig2005}
J.~J.~Craig, \textit{Introduction to Robotics: Mechanics and Control},
3rd~ed. Upper Saddle River, NJ: Pearson Prentice Hall, 2005.
ISBN: 978-0-13-123629-8.

\bibitem{cadena2016}
C.~Cadena \textit{et al.}, ``Past, present, and future of simultaneous
localization and mapping: Toward the robust-perception age,''
\textit{IEEE Trans.~Robotics}, vol.~32, no.~6, pp.~1309--1332, 2016.
\url{https://doi.org/10.1109/TRO.2016.2624754}

\end{thebibliography}

\end{document}
